%%% --- SYNTHÈSE DU CAS DU TÉLÉGRAPHE --- %%%

Le cas du télégraphe a fait émerger un ensemble d'observations dont certaines sont fermement établies par la littérature existante et d'autres restent au stade de l'hypothèse. Parmi les résultats les plus solides, on retiendra la substitution universelle du capital au travail~--- le même fait technique, documenté par \textcite{gordon_capital_nineteenth}, \textcite{holzmann_early_1995} et \textcite{gabler_american_1988}, se reproduit dans quatre contextes institutionnels radicalement différents~---, la convergence structurelle de ces quatre modèles vers le monopole~--- que \textcite{wu2010master} formalise et que \textcite{nonnenmacher_telegraph_2001}, \textcite{kieve_electric_1973} et \textcite{bertho_histoire_1984} documentent indépendamment~---, l'effet constitutif de l'information synchrone sur la formation des marchés~--- établi économétriquement par \textcite{steinwender_real_2018} et précisé par \textcite{juhasz_spinning_2018} pour les biens codifiables~---, le pattern récurrent d'investissement massif, de faillite et de recapitalisation qui accompagne le financement des câbles sous-marins~--- documenté par \textcite{gordon_thread_2002}, \textcite{bata_cables_1982} et \textcite{headrick_invisible_1991}~---, et enfin la matérialité occultée du réseau, massive mais invisible dans les comptes d'exploitation~--- documentée par \textcite{tully_victorian_2009}, \textcite{fitzmaurice_material_2023} et \textcite{hunt_imperial_2021}. Au stade de la piste, en revanche, on notera l'asymétrie entre la rivalité de la ligne et la non-rivalité de l'information transmise, observée mais non démontrée, et le rôle de la codifiabilité du produit dans la création de marchés, suggéré par \textcite{juhasz_spinning_2018} mais non testé systématiquement.

Ces observations ne sont pas indépendantes. La substitution du capital au travail détermine la structure de coûts~--- des coûts fixes élevés, des coûts marginaux faibles~--- qui pousse au monopole, quelle que soit la configuration institutionnelle de départ. Le monopole est à son tour la condition qui rend possible le déploiement de l'infrastructure à grande échelle~--- le câble transatlantique, le réseau national~---, et c'est cette infrastructure qui produit l'effet de coordination marchande documenté par Steinwender. La matérialité de cette infrastructure~--- gutta-percha, cuivre, bois~--- est occultée par l'attention que les contemporains portent aux prix et aux volumes du commerce qu'elle permet. Le circuit va de la structure du capital aux rendements croissants, des rendements croissants à la coordination marchande, et la matérialité de l'ensemble reste dans l'ombre.

Le passage d'un stade de maturité de la technologie au suivant n'a rien d'automatique, et le cas du télégraphe en fournit deux illustrations nettes. En France, l'administration du sémaphore optique bloqua le passage au télégraphe électrique pendant douze ans~--- de 1839 à 1851~--- par une combinaison de monopole d'État et de lobby institutionnel. En Grande-Bretagne, le \emph{block system} ferroviaire resta au stade du produit pendant quarante-sept ans~--- de 1842 à 1889~--- parce que le coût de l'installation était privé mais le bénéfice de la sécurité, public. Ce n'était pas, dans un cas comme dans l'autre, un retard d'adoption~: c'était une impossibilité structurelle de passage tant que les conditions~--- institutionnelles en France, économiques en Angleterre~--- n'étaient pas réunies.

Ce que le cas du télégraphe laisse ouvert, et que les cas suivants devront tester, tient en trois questions. Premièrement, l'asymétrie entre rivalité de l'infrastructure et non-rivalité de l'information transmise est-elle spécifique au télégraphe ou se retrouve-t-elle dans les technologies ultérieures, et comment se transforme-t-elle~? Le téléphone, qui occupe la ligne pendant toute la durée de la conversation, la pose différemment. Deuxièmement, le circuit causal qui va de la structure du capital au monopole, puis du monopole à la coordination marchande, se reproduit-il~--- et si oui, avec quelles variantes~? Le téléphone (AT\&T), la radio (RCA) et Internet (les plateformes) sont les candidats naturels. Troisièmement, les conditions de blocage observées en France et en Grande-Bretagne sont-elles des accidents locaux ou des manifestations d'un mécanisme récurrent~? Si le même type de blocage~--- un bénéfice public financé par un coût privé, un \emph{lobby} de l'ancienne technologie~--- se retrouve dans d'autres cas, il faudra en rendre compte comme d'une propriété structurelle, pas comme d'une anecdote.
